\section*{\S~8 Umgang mit Regelverstößen}
\label{sec:umgang_mit_regelverstoeßen}
\addcontentsline{toc}{section}{\S~8 Umgang mit Regelverstößen}

\begin{enumerate}[label=(\arabic*)]
    \item Jeder Regelverstoß wird geahndet.    
    \item Jedes Mitglied ist verpflichtet, grobe Regelverstöße zu melden.  
    \item\label{itm:vorgehen_bei_verstoß} Ein Regelverstoß zieht eine mündliche Verwarnung durch die Leitung nach sich. Dieser gilt für 3 Monate. Verstößt das verwarnte Mitglied in dieser Zeit gegen eine Regel, greift das nachfolgende Prozedere:
    \begin{enumerate}[label=(\Alph*)]
        \item Handelt es sich um einen Verstoß, der mindestens ein anderes Mitglied betrifft, lädt der Stab zu einem klärenden Gespräch mit allen Seiten ein. Der Gefreite vom Dienst (GvD) vertritt dabei die restliche Mannschaft.
        \item Handelt es sich um einen Verstoß, der nur die Person desjenigen betrifft, die den Regelverstoß begangen hat, lädt der Stab zu einem klärenden Gespräch mit dieser Person ein. Der Gefreite vom Dienst (GvD) kann hierbei die restliche Mannschaft vertreten.
    \end{enumerate}
    Kommt es trotz diesem Gespräch zu einem weiteren Regelverstoß, oder das Gespräch führt zu keiner Einigung, legt der Stab, nach Anhörung aller Seiten, eine Strafe fest. Der Gefreite vom Dienst (GvD) vertritt dabei in jedem Fall die restliche Mannschaft. Wird diese Entscheidung innerhalb von 72 Stunden nach Bekanntgabe von mindestens einem Mitglied angefochten, so tagt der erweiterte Stab nach erneuter Anhörung aller Seiten und legt eine Strafe fest. Diese Entscheidung gilt dann endgültig. Gibt es keinen erweiterten Stab, so tagt der Stab erneut. Eine Entscheidung muss in jedem Fall öffentlich bekannt gemacht werden.    
    \item\label{itm:schwerer_verstoß}Die Strafen richten sich nach der Schwere des Vergehens. Alle Vergehen, die das öffentliche Ansehen der 1st Aircat Division schaden und/oder andere Spieler im Spielspaß einschränken, gelten als schweres Vergehen.    
    \item Leichte Vergehen werden mit einem Ausschluss von einer Mission, oder einem Ausschluss von Beförderungen für einen Monat bestraft.    
    \item\label{itm:strafe_schwerer_verstoß} Schwere Vergehen werden mit einem Ausschluss von mehreren Missionen, mindestens aber 4 Missionen, oder einem Ausschluss von Beförderung für mehrere Monate, mindestens aber 3 Monate, bestraft.   
    \item Wiederholte schwere Vergehen werden mit einem Ausschluss von mehreren Missionen, mindestens aber 6 Missionen, und/oder einem Ausschluss von Beförderung für mehrere Monate, mindestens aber 6 Monaten, und/oder einer Degradierung bestraft. In besonders schweren Fällen kann das bestrafte Mitglied aus der Unit ausgeschlossen werden. Die Dauer des Ausschlusses legt das bestrafende Organ fest. Die Dauer des Ausschlusses darf nicht unter 6 Monaten liegen. Ein permanenter Ausschluss ist möglich.
\end{enumerate}
